\ifdefined\japsubmission
\begin{abstract}
Stochastic reaction networks model populations changed by random reaction
events.  A central stability question is whether such a system returns often
enough to admit a stationary probability distribution rather than escaping
permanently toward ever-larger populations.  We prove that every closed
communicating class has this stability, for every choice of positive reaction
rates, when the network is weakly reversible, has one linkage class, and no
reaction complex contains more than two molecules.  This removes an
additional pure-species condition required by the previous theorem.  The
proof follows the target produced by each reaction and measures the population
remaining after that target is subtracted.  Multiple linkage classes and
reactions involving more than two molecules remain open.
\end{abstract}
\else
\begin{abstract}
Weak reversibility is conjectured to imply positive recurrence for stochastic
mass-action systems.  For binary networks with one linkage class, this was
previously known under the additional requirement that every species occur in
a pure unary or pure-double complex.  We remove that requirement.  For every
positive rate vector, the minimal continuous-time Markov chain restricted to
each closed communicating class is nonexplosive.  Every nonabsorbing class is
positive recurrent, while an absorbing singleton carries its point-mass law;
hence every closed class admits a unique stationary probability distribution.
The proof marks
the target of the most recently fired reaction channel and applies a
log-factorial potential to the residual population obtained by subtracting
that target.  Its increment is exactly a target/source falling-factorial
ratio, so following the carried target has zero reward.  Finite directed paths
propagate the negative drift of a rare terminal source.  A normalized-log
compactification and an exhaustive bimolecular top-complex alternative either
supply such a source or produce a stoichiometric invariant that precludes
divergence within the class.  The result is class-wise and does not address
multiple linkage classes or molecularity above two.
\end{abstract}
\fi

\noindent\textbf{Keywords.}
Stochastic reaction network; systems biology; chemical master equation;
positive recurrence; mass-action kinetics; weak reversibility; continuous-time
Markov chain; Foster--Lyapunov method.

\medskip
\noindent\textbf{2020 Mathematics Subject Classification.}
Primary 60J28; secondary 60J27, 60J74, 92C42, 37N25.

\section{Introduction}\label{sec:intro}

Stochastic reaction networks are continuous-time Markov-chain models for
interacting populations whose transition rates are determined by a finite
reaction graph.  In systems biology they describe random molecule counts in
low-copy-number intracellular networks, including signalling and enzymatic
systems, gene-regulatory circuits, and synthetic biochemical networks
\citep{AndersonKurtz2015}.  Such models are often called chemical-master-equation
models.  For a nonexplosive nonabsorbing irreducible component, a basic
stability question
is whether it has a stationary probability distribution, or equivalently
whether its states have finite mean return times.  A stationary law gives a
long-run molecule-count distribution and expectations of bounded observables
within that component; moments of unbounded count observables require separate
integrability.  Positive recurrence does not imply that sample paths remain
in a bounded set.

The deterministic counterpart supplies additional motivation.  Weakly
reversible single-linkage mass-action systems are bounded
\citep{Anderson2011Boundedness} and permanent \citep{BorosHofbauer2020} under
their deterministic dynamics.  Those results express a strong structural
stability of concentrations, but deterministic boundedness or permanence does
not by itself prove recurrence of the discrete stochastic model.  Positive
recurrence is the corresponding class-wise stochastic question addressed
here.

\citet{AndersonKim2018} formulated the positive-recurrence conjecture that
weak reversibility should force this stability for stochastic mass-action
systems.  They established several structural sufficient conditions.
\citet{AndersonCappellettiKim2020} subsequently proved the binary,
one-linkage case under the additional condition
\begin{equation}\label{eq:ACK-condition}
  \{S_i,2S_i\}\cap\mathcal C\ne\varnothing
  \qquad\text{for every species }S_i.
\end{equation}
Thus each species had to occur in a pure unary or pure double complex.  The
proof of \citet{AndersonCappellettiKim2020} applies an entropy-like
function to an \(n\)-step embedded chain.  In their Section~6, Theorem~6.1
reduces recurrence to their tier inclusion (11): every top S-tier source must
also lie in the top D-tier along every proper tier sequence.  Lemmas~6.4--6.5
then construct a bounded reaction word and convert its D-tier drop into
negative sampled-chain drift.  The pure-species assumption enters decisively
in the proof of their Theorem~4.1: for a D-top mixed complex \(S_u+S_v\) that
is disabled because the \(u\)-coordinate is zero while the \(v\)-coordinate
diverges, the existence of \(S_v\) or \(2S_v\) supplies the comparison needed
to establish (11) \citep[Section~6.1, equations~(19)--(20), pp.~808--809]
{AndersonCappellettiKim2020}.  The marked-target construction used here avoids
that boundary-availability step: the carried target is automatically enabled
at the post-jump state, and a directed target-following path propagates its
reward until the top-complex alternative supplies a rare terminal source.  It
therefore does not require a unary or pure-double complex for every species.

Other progress includes complex-balanced product-form systems
\citep{AndersonCraciunKurtz2010}, strongly endotactic structures
\citep{AndersonCappellettiKimNguyen2020}, one-dimensional classifications
\citep{WiufXu2023}, a path method for exponential ergodicity
\citep{AndersonCappellettiFanKim2025}, and nonexplosion of all second-order
weakly reversible systems \citep{Xu2026}.  A complementary two-species
positive-recurrence theorem has been publicly announced in conference talks
since 2022 \citep{Agazzi2022Talk,Agazzi2025Talk}.  The current ConStRAINeD
project page lists a five-author proof as in preparation
\citep{ConStRAINeDResults}.  We did not locate a public manuscript for that
result as of 9 August 2026.  The public descriptions concern two species and
permit broader network structure; the present theorem permits arbitrarily many
species but assumes one linkage class and molecularity at most two.  Neither
result should be presented as superseding the other without the announced
manuscript.  The May 2026
revision of \citet{Xu2026} still records the bimolecular positive-recurrence
problem as open.

The present paper removes \eqref{eq:ACK-condition} while retaining the binary
and one-linkage assumptions; in particular, its theorem contains the result of
\citet{AndersonCappellettiKim2020} as a special case.  The exact theorem is
stated in Theorem~\ref{thm:main}.  It concerns each already-closed
communicating class.  It does not prove finite expected entrance, from an
arbitrary nonclosed initial state, into the union of closed irreducible
components.  Nor does it address multiple linkage classes or complexes of
molecularity above two, and therefore it does not settle the full
positive-recurrence conjecture.

A small motivating example is the directed cycle
\begin{equation}\label{eq:stress}
  0\longrightarrow A+B\longrightarrow B\longrightarrow0.
\end{equation}
Species \(A\) occurs in neither \(A\) nor \(2A\), so
\eqref{eq:ACK-condition} fails.  At the boundary state \((n,0)\), only the
reaction \(0\to A+B\) is enabled and the one-step drift of total population
is positive.  The restoring effect is delayed: the created \(B\) enables a
fast reaction consuming \(A\).  This is representative of why a one-step
radial Lyapunov function is not the right universal object.

The delayed correction can be seen directly.  Give the three channels in
\eqref{eq:stress} rate constants \(\kappa_0,\kappa_1,\kappa_2\), in the order
displayed, and begin at the reachable marked state \(((n,0),0)\).  The two
designated target-following reactions \(0\to A+B\) and \(A+B\to B\) each
have zero residual-potential increment.  The second is selected with
probability
\[
 \frac{\kappa_1(n+1)}{\kappa_0+\kappa_2+\kappa_1(n+1)}=1-O(n^{-1}).
\]
At the resulting marked state \(((n,1),B)\), the expected next increment is
\begin{equation}\label{eq:stress-drift}
 -\frac{\kappa_1n}{\kappa_0+\kappa_2+\kappa_1n}\log n
 =-\log n+o(\log n).
\end{equation}
The exceptional branches at the preceding phase have probability
\(O(n^{-1})\) and increments of order at most \(\log n\).  Thus the complete
target-following episode has a restoring drift of order \(-\log n\), even
though the first boundary jump increases the total population.

The proof instead uses four ideas.  First, we augment the embedded jump chain
by marking the actual reaction channel that fired and retaining its target
complex.  Second, after removing that target from the population, the
log-factorial potential has the exact increment
\[
  \log\frac{(x)_t}{(x)_s}
\]
when the carried target is \(t\) and the next source is \(s\).  Following the
carried target therefore costs exactly zero.  Third, a finite episode follows
a fixed directed path from the carried target toward a selected terminal
complex.  A sharp scalar envelope shows that a terminal source probability
tending to zero forces the complete episode drift to tend to \(-\infty\), no
matter how the intermediate source rates are separated.  Fourth, a
normalized-log compactification of an arbitrary divergent sequence gives an
exhaustive top-complex case split.  Bimolecularity implies either the
existence of a higher-weight source enabled over a lower terminal complex or
a reaction-wise linear invariant that contradicts divergence inside a fixed
communicating class.  A finite-family random-time drift argument and a finite
trace-chain construction complete the proof.

\section{Reaction networks and the stochastic model}\label{sec:model}

Let \(\mathcal S=\{S_1,\dots,S_d\}\) be a finite species set.  A
\emph{complex} is a vector \(y=(y_1,\dots,y_d)\in\mathbb N_0^d\), identified
with the formal sum \(\sum_i y_iS_i\).  Its molecularity is
\(|y|=\sum_i y_i\).  A \emph{reaction channel} is a labelled ordered pair
\(r:y_r\to y_r'\), together with a positive rate constant \(\kappa_r\).
Parallel channels are allowed.  The reaction graph has the complexes as
vertices and the reaction channels as directed edges.  Its linkage classes
are the connected components of the underlying undirected graph.  A network
is \emph{weakly reversible} when each linkage class is strongly connected.
With one linkage class, weak reversibility is simply strong connectivity of
the entire complex graph.

We call the network \emph{binary} or \emph{bimolecular} when
\begin{equation}\label{eq:binary}
  |y|\le2\qquad\text{for every }y\in\mathcal C.
\end{equation}
Because every complex in a weakly reversible graph is a reaction source, this
agrees here with the usual second-order/bimolecular terminology.

For \(x\in\mathbb N_0^d\), define the multivariate falling factorial
\begin{equation}\label{eq:falling}
  (x)_y=
  \begin{cases}
  \displaystyle\prod_{i=1}^d\frac{x_i!}{(x_i-y_i)!},&x\ge y,\\[1ex]
  0,&\text{otherwise}.
  \end{cases}
\end{equation}
The stochastic mass-action propensity of channel \(r\) is
\(\lambda_r(x)=\kappa_r(x)_{y_r}\), and the formal generator is
\begin{equation}\label{eq:generator}
  \mathcal Lf(x)=\sum_r\kappa_r(x)_{y_r}
  \bigl[f(x+y_r'-y_r)-f(x)\bigr].
\end{equation}
The convention using products of binomial coefficients is equivalent after
rescaling each positive rate constant by a fixed factorial factor.

A subset \(\Gamma\subseteq\mathbb N_0^d\) is a \emph{closed communicating
class} when its states communicate through enabled channels and no enabled
channel exits \(\Gamma\).  The chain restricted to \(\Gamma\) is
irreducible.  In a nonabsorbing class, let \(J_1\) be the first jump time and
define the positive return time to \(x\) by
\(T_x^+=\inf\{t\ge J_1:X(t)=x\}\).  The class is \emph{positive recurrent}
when one, equivalently every, state has \(\mathbb E_xT_x^+<\infty\).
Absorbing singleton classes are handled separately: their stationary
probability is the point mass at the absorbing state.  Henceforth positive
return times are discussed only for nonabsorbing irreducible classes.  For a
nonexplosive irreducible countable-state CTMC, positive recurrence is
equivalent to the existence of a stationary probability distribution
\citep{Norris1997}.

A channel \(y\to y\) has zero population increment and contributes zero to
\eqref{eq:generator}; discard it.  Exact parallel channels with the same
source and target may be combined by adding their rates.  Channels having the
same population displacement but different sources or targets remain
distinct, because the proof marks the actual channel.
All reaction-indexed sums and rate quantities below, including \(\Lambda\)
and \(\bar\kappa_y\), refer to this reduced set of genuine channels.

Every finite nonabsorbing closed communicating class is positive recurrent:
its finite irreducible jump chain has a stationary probability distribution,
all state rates are finite, and the associated finite-state CTMC is
nonexplosive \citep{Norris1997}.  If
the complex graph has one vertex, all closed classes are absorbing
singletons.  If all complexes have the same molecularity, every reaction
preserves \(|x|_1\), so each communicating class lies in a finite shell.  We
therefore fix henceforth an infinite closed irreducible class \(\Gamma\), a
strongly connected graph with at least two complexes, and no self channels.
Every complex then has a genuine outgoing channel.  Write
\begin{equation}\label{eq:kbar}
  \bar\kappa_y=\sum_{r:y_r=y}\kappa_r,
  \qquad
  \bar\kappa_- =\min_{y\in\mathcal C}\bar\kappa_y>0,
  \qquad
  \bar\kappa_+ =\max_{y\in\mathcal C}\bar\kappa_y.
\end{equation}

\begin{theorem}[Binary one-linkage positive recurrence]\label{thm:main}
Let a finite stochastic mass-action reaction network be weakly reversible,
have one linkage class, and satisfy \eqref{eq:binary}.  For every positive
rate vector, the minimal population CTMC restricted to each closed
communicating class is nonexplosive.  Every nonabsorbing such class is
positive recurrent; each absorbing singleton has its point-mass stationary
law.
\end{theorem}

\begin{corollary}[Stationary probability]\label{cor:stationary}
Under the hypotheses of Theorem~\ref{thm:main}, every closed communicating
class \(\Gamma\) admits a stationary probability distribution
\(\pi_\Gamma\).  It is unique on \(\Gamma\).
\end{corollary}

\section{The marked target-augmented chain}\label{sec:marked}

We first define the embedded chain without assuming nonexplosion.  At a
nonabsorbing population state \(x\), sample the actual labelled reaction
channel \(r\) with probability
\begin{equation}\label{eq:channel-kernel}
  \frac{\kappa_r(x)_{y_r}}{\Lambda(x)},
  \qquad
  \Lambda(x)=\sum_q\kappa_q(x)_{y_q}.
\end{equation}
This discrete chain can be constructed independently of its holding times.
After the sampled channel \(r\) fires, record its target \(t=y_r'\).  The mark
is therefore obtained from the actual channel and is never inferred from the
population displacement.  Once exact parallel channels with the same source
and target have been aggregated, retaining the target is the only part of the
previous channel needed below.  Write \(Z_n\) for this target-augmented
embedded chain.  A reachable augmented state is denoted \((x,t)\), where
\(x\) is the post-jump population.  If the next sampled
channel is \(q:s\to u\), then
\begin{equation}\label{eq:aug-transition}
  (x,t)\longmapsto(x-s+u,u).
\end{equation}
The channel law depends on \(x\), not on the historical mark, so the augmented
process is Markov, and projection onto \(x\) is the ordinary population
embedded chain.

Let \(\widetilde\Gamma\) be the augmented states reachable after genuine
jumps within \(\Gamma\).  It is closed under \eqref{eq:aug-transition}.  It is
also irreducible, and its projection onto the population coordinate is exactly
\(\Gamma\).  For the latter assertion, every \(x\in\Gamma\) terminates a
positive-length population path: irreducibility gives such a path directly,
or, if necessary, one may take a genuine outgoing jump and a return path.
The final channel supplies a reachable mark over \(x\).  For irreducibility,
given \((x,t),(x',t')\in\widetilde\Gamma\), choose
a channel \(s'\to t'\) and predecessor population
\(z=x'-t'+s'\in\Gamma\) witnessing reachability of the second pair.
Population irreducibility supplies a positive-probability marked-channel path
from \(x\) to \(z\); appending \(s'\to t'\) reaches \((x',t')\).

The carried target is always enabled.  If the preceding channel was
\(s\to t\), then its post-jump state has the form \(x=z-s+t\ge t\).  Put
\begin{equation}\label{eq:residual}
  r=x-t\in\mathbb N_0^d,
  \qquad
  V(x,t)=\sum_{i=1}^d\log((x_i-t_i)!)=\sum_i\log(r_i!).
\end{equation}
We use \(\log(0!)=0\).  Thus \(V\ge0\).  Every sublevel set of \(V\) in
\(\widetilde\Gamma\) is finite, because \(\log(n!)\to\infty\), each residual
coordinate is bounded on a sublevel, and the target set is finite.

The entropy growth in \eqref{eq:residual} is classical.  With
\(0\log0=0\), Stirling's formula gives
\begin{equation}\label{eq:stirling-entropy}
 \sum_i\log(r_i!)
 =\sum_i\bigl(r_i\log r_i-r_i\bigr)
  +O\!\left(\sum_i\log(r_i+1)\right).
\end{equation}
This leading term is in the pseudo-Helmholtz or Horn--Jackson family
\[
 G_c(r)=\sum_i\left[r_i\bigl(\log(r_i/c_i)-1\bigr)+c_i\right],
 \qquad c_i>0,
\]
used in deterministic and stochastic chemical reaction network theory
\citep{HornJackson1972,AndersonCraciunGopalkrishnanWiuf2015}.  Thus the residual log-factorial
potential is a discrete, target-shifted analogue of that classical entropy,
not the Horn--Jackson function itself.  The target shift is the new feature in
this argument: subtracting the complex actually produced by the preceding
channel gives the exact identity below, rather than only an asymptotic
Lyapunov estimate, and makes every target-following jump have exactly zero
increment.

\begin{lemma}[Exact target/source identity]\label{lem:identity}
At \((x,t)\), if the next marked channel has source \(s\) and target \(u\),
then
\begin{equation}\label{eq:identity}
  V(x-s+u,u)-V(x,t)=\log\frac{(x)_t}{(x)_s}.
\end{equation}
In particular, a channel sourced at the carried target has zero increment.
\end{lemma}

\begin{proof}
The new residual is \((x-s+u)-u=x-s\).  Since both \(s\) and \(t\) are
enabled,
\[
 \exp(\Delta V)
 =\prod_i\frac{(x_i-s_i)!}{(x_i-t_i)!}
 =\frac{(x)_t}{(x)_s}.
\]
\end{proof}

This identity gives a precise ``complete-credit'' interpretation: a product
complex remains credited in full until the chain chooses a different source.
The interpretation is not used in place of the algebraic identity.

For enabled source complexes define
\begin{equation}\label{eq:source-prob}
  p_x(y)=\frac{\bar\kappa_y(x)_y}{\Lambda(x)}.
\end{equation}
Set \(p_x(y)=0\) for disabled sources, and use \(0\log0=0\) in entropy
expressions.  Unless explicitly stated otherwise, sums in this subsection
range only over enabled source complexes.  Because the carried target is
enabled, \(p_x(t)>0\).  The expected one-jump increment is
\begin{equation}\label{eq:d-def}
  d(x,t)=\sum_{s:\,x\ge s}p_x(s)
  \log\frac{(x)_t}{(x)_s}.
\end{equation}

\begin{lemma}[Source-probability bound]\label{lem:entropy}
For every reachable augmented state,
\begin{equation}\label{eq:entropy-identity}
\begin{split}
  d(x,t)
  ={}&\log p_x(t)-\sum_{s:\,x\ge s} p_x(s)\log p_x(s)\\
    &+\sum_{s:\,x\ge s} p_x(s)\log\bar\kappa_s
      -\log\bar\kappa_t,
\end{split}
\end{equation}
and consequently
\begin{equation}\label{eq:entropy-bound}
  d(x,t)\le\log p_x(t)+C_0,
  \qquad
  C_0=\log|\mathcal C|+
  \log\frac{\bar\kappa_+}{\bar\kappa_-}.
\end{equation}
\end{lemma}

\begin{proof}
For each enabled \(s\),
\((x)_s=p_x(s)\Lambda(x)/\bar\kappa_s\).  Substitution into
\eqref{eq:d-def} gives \eqref{eq:entropy-identity}.  The entropy of a
probability vector supported on at most \(|\mathcal C|\) points is at most
\(\log|\mathcal C|\), and the last line contributes at most
\(\log(\bar\kappa_+/\bar\kappa_-)\).  The zero complex creates no exception,
since \((x)_0=1\).
\end{proof}

\section{Propagation along target-following paths}\label{sec:paths}

For every ordered pair \((t,c)\in\mathcal C^2\), fix a simple directed path
of marked channels
\begin{equation}\label{eq:path}
 t=y_0\longrightarrow y_1\longrightarrow\cdots
 \longrightarrow y_L=c,
 \qquad 0\le L\le|\mathcal C|-1.
\end{equation}
Strong connectivity supplies the path.  Starting from \((r+t,t)\), define
the \((t,c)\)-episode as follows.  At phase \(y_k\), \(k<L\), continue only
when the exact designated channel \(y_k\to y_{k+1}\) fires; after any other
channel, stop immediately.  If the path reaches \(c\), take one final
ordinary jump and stop.  When \(c=t\), the path has length zero and the
episode is just that final jump.

This is a stopping time for the marked embedded filtration.  On designated
edges the lifted populations are
\begin{equation}\label{eq:lift}
  r+y_0\longrightarrow r+y_1\longrightarrow\cdots\longrightarrow r+c.
\end{equation}
Every designated source is literally present, the residual remains \(r\),
and closedness of \(\Gamma\) keeps the entire lift in the class.  The episode
has between one and \(|\mathcal C|\) jumps and only finitely many terminal
branches.

For the designated channel at phase \(k\), put
\begin{equation}\label{eq:qk}
 q_k=\frac{\kappa_{y_k\to y_{k+1}}}{\bar\kappa_{y_k}}>0,
 \qquad p_k=p_{r+y_k}(y_k).
\end{equation}
Let \(J_k(r)\) be the expected remaining \(V\)-increment from phase \(k\).

\begin{lemma}[Exact episode recursion]\label{lem:recursion}
The episode rewards satisfy
\begin{equation}\label{eq:recursion}
  J_L(r)=d(r+c,c),
  \qquad
  J_k(r)=d(r+y_k,y_k)+q_kp_kJ_{k+1}(r).
\end{equation}
\end{lemma}

\begin{proof}
At phase \(k\), the first jump has expected increment
\(d(r+y_k,y_k)\).  Continuation occurs exactly when the designated channel
fires, with probability \(q_kp_k\).  That channel is sourced at the carried
target and has zero immediate increment by Lemma~\ref{lem:identity}; the
remaining conditional expectation is \(J_{k+1}(r)\).
\end{proof}

The following scalar calculation controls all possible intermediate source
probabilities.

\begin{lemma}[Scalar envelope]\label{lem:envelope}
For \(q>0\) and \(M\in\mathbb R\), define
\[
  F_q(M)=\sup_{0<p\le1}\{\log p+C_0+qpM\}.
\]
The function \(F_q\) is nondecreasing in \(M\).
Then
\begin{equation}\label{eq:F-piecewise}
  F_q(M)=
  \begin{cases}
    C_0+qM,&M\ge-1/q,\\
    C_0-1-\log(-qM),&M<-1/q.
  \end{cases}
\end{equation}
Consequently, if \(M_n\to-\infty\), then \(F_q(M_n)\to-\infty\).
\end{lemma}

\begin{proof}
For each fixed \(p\in(0,1]\), the displayed objective is increasing in
\(M\), with slope \(qp>0\); taking its supremum proves monotonicity.
The objective has derivative \(p^{-1}+qM\) and second derivative
\(-p^{-2}<0\).  If \(M\ge-1/q\), its derivative at \(p=1\) is nonnegative,
so concavity places the maximum at \(p=1\).  If \(M<-1/q\), the unique
maximizer is \(p=(-qM)^{-1}\in(0,1)\), which gives the second line of
\eqref{eq:F-piecewise}.  The final assertion follows immediately.
\end{proof}

\begin{proposition}[Propagation of terminal drift]\label{prop:propagation}
Fix a target-following path \eqref{eq:path}.  If along a sequence of residuals
\(r^{(n)}\) one has
\[
  p_{r^{(n)}+c}(c)\longrightarrow0,
\]
then the expected increment of the complete \((t,c)\)-episode tends to
\(-\infty\).
\end{proposition}

\begin{proof}
At the terminal phase, Lemma~\ref{lem:entropy} gives
\(J_L\le\log p_{r+c}(c)+C_0\to-\infty\).  If \(J_{k+1}\le M\), then the
monotonicity just proved and
Lemmas~\ref{lem:entropy} and \ref{lem:recursion} give
\[
  J_k\le\sup_{0<p\le1}\{\log p+C_0+q_kpM\}=F_{q_k}(M).
\]
Apply Lemma~\ref{lem:envelope} backward through the finite path.
\end{proof}

The proof preserves every reaction-rate monomial: no edge is replaced and no
unrelated rates are compared.  Arbitrarily many nested source-rate scales are
absorbed by the finite scalar recursion.

\section{Logarithmic compactification and the top-complex alternative}
\label{sec:compact}

We now show that every divergent augmented sequence supplies the terminal
rarity required by Proposition~\ref{prop:propagation}, unless a linear
stoichiometric invariant already makes such divergence impossible.

Take a sequence \((x^{(n)},t^{(n)})\in\widetilde\Gamma\) leaving every finite
set.  Pass to a subsequence with \(t^{(n)}=t\), and set
\(r^{(n)}=x^{(n)}-t\).  A diagonal extraction makes every coordinate of
\(r^{(n)}\) either one fixed nonnegative integer or divergent to infinity.
Let \(I\) be the set of divergent coordinates and put
\begin{equation}\label{eq:Rn}
  R_n=\sum_{i\in I}\log(r_i^{(n)}+1).
\end{equation}
After a further subsequence, the limits
\begin{equation}\label{eq:weights}
  w_i=\lim_{n\to\infty}
  \frac{\log(r_i^{(n)}+1)}{R_n}
  \quad(i\in I),
  \qquad w_i=0\quad(i\notin I)
\end{equation}
exist.  Thus \(w\ge0\) and \(\sum_iw_i=1\).

\begin{remark}[Slower divergent tiers]\label{rem:zero-weight}
A coordinate can diverge while having \(w_i=0\).  It remains in \(I\): zero
normalized logarithmic weight means growth on a slower divergent tier, not a
bounded species.  This distinction is essential in the availability and
signed-invariant branches of Lemma~\ref{lem:top}.
\end{remark}

\begin{lemma}[Falling-factorial asymptotics]\label{lem:ff-asymptotic}
Let \(c\) and \(y\) be fixed complexes, and suppose \(y\) is enabled at
\(r^{(n)}+c\) for all sufficiently large \(n\).  Then
\begin{equation}\label{eq:ff-asymptotic}
  \log(r^{(n)}+c)_y=R_n\,w\cdot y+o(R_n).
\end{equation}
\end{lemma}

\begin{proof}
For \(y_i=1\) and a divergent coordinate,
\(\log(r_i^{(n)}+c_i)=\log(r_i^{(n)}+1)+o(R_n)\); a fixed coordinate
contributes \(O(1)\).  For \(y_i=2\) and \(r_i^{(n)}\to\infty\),
\[
 \log\bigl((r_i^{(n)}+c_i)(r_i^{(n)}+c_i-1)\bigr)
 =2\log(r_i^{(n)}+1)+o(R_n).
\]
If that coordinate is fixed, enabledness makes its contribution finite and
hence \(O(1)\).  Summing the coordinate contributions proves the claim.
\end{proof}

For a complex \(y\), write \(h(y)=w\cdot y\), and define
\begin{equation}\label{eq:top}
  a=\max_{y\in\mathcal C}h(y),
  \qquad
  T=\{y\in\mathcal C:h(y)=a\}.
\end{equation}
The next lemma contains the decisive use of bimolecularity.
When \(T\ne\mathcal C\), its maximal weight satisfies \(a>0\), so every top
complex contains at least one particle in a divergent coordinate.  Because a
binary complex contains at most two particles, only two configurations remain:
a top complex contains two such particles, or every top complex contains
exactly one.  This elementary restriction makes the following alternative
exhaustive; it is precisely what fails without the molecularity bound.

\begin{lemma}[Top-complex availability or invariant]\label{lem:top}
Along the divergent sequence, at least one of the following conclusions is
available:

\begin{enumerate}[label=\textup{(\roman*)},leftmargin=2.5em]
\item a reaction-wise linear stoichiometric invariant has value tending to
      infinity along the sequence, contradicting membership in the fixed
      communicating class; or
\item there are fixed complexes \(s,c\in\mathcal C\) with
      \(h(s)>h(c)\) such that both \(s\) and \(c\) are enabled at
      \(r^{(n)}+c\) for all sufficiently large \(n\).
\end{enumerate}
\end{lemma}

\begin{proof}
We follow the three exhaustive cases described above.  First, if every
complex is top, then \(w\cdot(y'-y)=0\) for every channel.
Therefore \(w\cdot X\) is a nonnegative linear invariant.  A coordinate with
positive weight diverges, so the invariant value tends to infinity, which is
impossible inside one communicating class.

Second, assume \(T\ne\mathcal C\) and a top complex \(s\) contains two
particles supported on coordinates in \(I\).  Bimolecularity forces \(s\) to
contain no particle supported on a bounded coordinate.  Both required
particles are eventually present in \(r^{(n)}\), including the case
\(s=2S_i\).  Hence \(s\) is enabled at \(r^{(n)}+c\) for every lower
\(c\notin T\), and \(h(s)>h(c)\).

Third, suppose every top complex contains exactly one particle supported on a
coordinate in \(I\).  Let \(J\subseteq I\) be the species appearing in top
complexes.  If \(i\in J\), the remaining possible particle in that top
complex is supported on a bounded coordinate and therefore has weight zero;
consequently \(w_i=a\).  No complex contains two \(J\)-particles, since its
weight would be at least \(2a>a\).  Moreover,
\begin{equation}\label{eq:top-iff}
  y\in T\quad\Longleftrightarrow\quad
  q_J(y):=\sum_{i\in J}y_i=1.
\end{equation}
The forward implication follows from the definition of \(J\).  Conversely,
one \(J\)-particle already contributes \(a\); maximality forces all other
weight to be zero and makes the complex top.

If every complex had \(q_J=1\), then \eqref{eq:top-iff} would give
\(T=\mathcal C\), the all-top case already handled above.  There remain three
subcases.

\begin{enumerate}[label=\textup{(\alph*)},leftmargin=2.5em]
\item If a unary top complex \(S_i\), \(i\in J\), exists, it is enabled over
      every lower terminal complex, giving conclusion (ii).
\item Otherwise every top complex has the form \(S_i+D\), where
      \(D\notin I\); call \(D\) a bounded companion species.  If a lower complex
      \(c\) contains such a \(D\), the corresponding top source \(S_i+D\)
      is enabled at \(r^{(n)}+c\), giving conclusion (ii).
\item If no lower complex contains any bounded companion species, let
      \(\mathcal D\) be the set of distinct companion species.  The signed
      linear functional
      \begin{equation}\label{eq:signed-invariant}
        \ell(x):=\sum_{i\in J}x_i-\sum_{D\in\mathcal D}x_D
      \end{equation}
      has value zero on every complex: each top complex contains one
      \(J\)-particle and one companion particle, whereas a lower complex
      contains neither.  It is therefore reaction-wise invariant.  Its
      positive part diverges, while every companion coordinate lies outside
      \(I\) and is fixed along the subsequence.  Its value tends to infinity,
      contradicting the fixed communicating class.
\end{enumerate}

This exhausts all binary complexes.  We reserve ``conservation law'' for the
nonnegative invariants above; \eqref{eq:signed-invariant} is a signed linear
stoichiometric invariant.
\end{proof}

The invariant contradiction is exact: if \(\ell\cdot(y'-y)=0\) on every
channel, then \(\ell\cdot x\) is constant along every path in \(\Gamma\).
In the signed branch, the coordinates with negative coefficients are fixed along
the extracted sequence, so the invariant really does tend to \(+\infty\).
A species absent from every complex is itself reaction-wise invariant; thus an
absent coordinate cannot diverge along a fixed communicating class.

\begin{lemma}[Vanishing terminal source probability]\label{lem:terminal}
Under conclusion \textup{(ii)} of Lemma~\ref{lem:top},
\begin{equation}\label{eq:terminal-zero}
  p_{r^{(n)}+c}(c)\longrightarrow0.
\end{equation}
\end{lemma}

\begin{proof}
Both \(s\) and \(c\) are enabled.  Lemma~\ref{lem:ff-asymptotic} and the
strict gap \(h(s)>h(c)\) give
\[
  \frac{(r^{(n)}+c)_s}{(r^{(n)}+c)_c}\longrightarrow\infty.
\]
Consequently,
\[
 p_{r^{(n)}+c}(c)
 \le
 \frac{\bar\kappa_c(r^{(n)}+c)_c}
      {\bar\kappa_s(r^{(n)}+c)_s}
 \longrightarrow0.
\]
\end{proof}

\section{Uniform random-time drift outside a finite set}\label{sec:uniform}

For each \(c\in\mathcal C\), let \(D_c(x,t)\) denote the exact expected
increment of the fixed \((t,c)\)-episode.  Define
\begin{equation}\label{eq:K}
 K=\left\{(x,t)\in\widetilde\Gamma:
          \min_{c\in\mathcal C}D_c(x,t)>-1\right\}.
\end{equation}

\begin{proposition}[Finite and nonempty exceptional set]\label{prop:K}
The set \(K\) is finite and nonempty.
\end{proposition}

\begin{proof}
Suppose first that \(K\) were infinite.  Properness of \(V\) supplies a
sequence in \(K\) leaving every finite set.  Apply the compactification of
Section~\ref{sec:compact}.  An invariant conclusion from
Lemma~\ref{lem:top} contradicts membership in one communicating class.
Otherwise, Lemma~\ref{lem:terminal} and
Proposition~\ref{prop:propagation} supply a fixed terminal complex \(c\) for
which \(D_c(x^{(n)},t)\to-\infty\).  This contradicts the definition of
\(K\).  Thus \(K\) is finite.

To prove nonemptiness, select any augmented state \(z_0\).  The finite
nonempty sublevel set \(\{z:V(z)\le V(z_0)\}\) contains a global minimizer
\(z_*\) of \(V\).  Every endpoint of every target-following episode remains
in \(\widetilde\Gamma\), so its potential is at least \(V(z_*)\).  Hence
\(D_c(z_*)\ge0\) for every \(c\), and \(z_*\in K\).
\end{proof}

\begin{remark}[Rate dependence of the exceptional set]\label{rem:rate-dependence}
The finiteness proof is qualitative and cannot give a rate-independent bound
on the location of \(K\).  Consider the directed cycle
\[
 0\xrightarrow{\kappa_0}A
 \xrightarrow{\kappa_1}A+B
 \xrightarrow{\kappa_2}0
\]
and the augmented state \(((m,0),A)\), \(m\ge2\).  The population chain is
irreducible on \(\mathbb N_0^2\), and each such marked state is reachable
after a \(0\to A\) jump from \((m-1,0)\); hence all these states lie in one
fixed augmented class.  For the target-following path from \(A\) through
\(A+B\) to terminal \(0\), the exact recursion gives
\begin{equation}\label{eq:rate-example}
 D_0(m,A)=a_m+p_m\bigl(b_m+q_mc_m\bigr),
\end{equation}
where
\[
\begin{gathered}
 a_m=\frac{\kappa_0\log m}{\kappa_0+\kappa_1m},\qquad
 p_m=\frac{\kappa_1m}{\kappa_0+\kappa_1m},\\
 b_m=\frac{\kappa_0\log m}{\kappa_0+(\kappa_1+\kappa_2)m},\qquad
 q_m=\frac{\kappa_2m}{\kappa_0+(\kappa_1+\kappa_2)m},\\
 c_m=-\frac{\kappa_1(m-1)}{\kappa_0+\kappa_1(m-1)}\log(m-1).
\end{gathered}
\]
The three terms are respectively the drift at carried targets \(A\),
\(A+B\), and \(0\), with \(p_m\) and \(q_m\) the exact continuation
probabilities.  Hence
\begin{equation}\label{eq:rate-example-asymptotic}
 D_0(m,A)=-\frac{\kappa_2}{\kappa_1+\kappa_2}\log m
 +O\!\left(\frac{\log m}{m}\right).
\end{equation}
For the other two possible terminals, the same path library gives
\(D_A(m,A)=a_m>0\) and
\(D_{A+B}(m,A)=a_m+p_mb_m>0\).  With \(\kappa_0,\kappa_1\) fixed,
\eqref{eq:rate-example} tends to \(a_m+p_mb_m>0\) as
\(\kappa_2\downarrow0\) for each fixed \(m\).  Thus, for every finite
\(M\), a sufficiently small positive \(\kappa_2\) puts all the states
\(((m,0),A)\), \(2\le m\le M\), in \(K\).  At the same time the negative
coefficient in \eqref{eq:rate-example-asymptotic} can be arbitrarily small.
Consequently, neither the location nor the diameter of \(K\) is bounded by the
numbers of species and complexes alone, uniformly over positive rate vectors.
\end{remark}

Outside \(K\), choose the lexicographically first terminal complex minimizing
\(D_c(x,t)\), before the episode begins.  This is a deterministic measurable
selector on the countable state space and is compatible with restarting
under the strong Markov property.  Its episode stopping time \(\tau\) obeys
\begin{equation}\label{eq:episode-drift}
  1\le\tau\le L:=|\mathcal C|,
  \qquad
  \mathbb E_{(x,t)}[V(Z_\tau)-V(x,t)]\le-1,
  \quad (x,t)\notin K.
\end{equation}
The coordinate overshoot in one episode is deterministically bounded by
\(2L\), because each binary reaction changes each population coordinate by
at most two.

\section{Positive recurrence}\label{sec:recurrence}

The following argument is a short specialization of standard sampled-chain
Foster theory \citep{MeynTweedie2009}; we include it to record the exact
stopping and return interfaces used here.

Let \(Y_0=z\).  If \(Y_n\notin K\), let \(Y_{n+1}\) be the endpoint of the
episode selected at \(Y_n\); if \(Y_n\in K\), set \(Y_{n+1}=Y_n\).  Thus the
endpoint chain is stopped upon entering \(K\).  Put
\(
 \sigma_K=\inf\{n\ge0:Y_n\in K\}.
\)
Let \(\mathcal G_n\) be the endpoint-chain filtration and define
\[
 W_n=V(Y_{n\wedge\sigma_K})+(n\wedge\sigma_K).
\]
The drift bound \eqref{eq:episode-drift} makes \((W_n)\) a nonnegative
supermartingale.  Its terms are integrable: after at most \(N\) episodes, each
population coordinate is bounded above by its initial value plus \(2LN\), by
the deterministic overshoot bound following \eqref{eq:episode-drift}; the
finite target set then bounds \(V\).  Applying the supermartingale inequality
at the bounded time \(N\) gives
\begin{equation}\label{eq:stopped-drift}
  \mathbb E_zV(Y_{N\wedge\sigma_K})
  +\mathbb E_z(N\wedge\sigma_K)
  \le V(z).
\end{equation}
Because \(V\ge0\), expectations and monotone convergence yield
\begin{equation}\label{eq:episodes-to-K}
  \mathbb E_z\sigma_K\le V(z).
\end{equation}
The original augmented embedded chain therefore hits \(K\) in finite
expected jump count, bounded by \(L V(z)\).

The following finite trace argument turns hitting \(K\) into positive return
to one state; the finite-chain step is standard \citep{Norris1997}, while the
conversion from trace excursions to original jumps is recorded explicitly.

\begin{proposition}[Finite trace-chain closure]\label{prop:trace}
The augmented embedded chain has a state with finite expected positive return
time.
\end{proposition}

\begin{proof}
Let \(T_K^+=\inf\{n\ge1:Z_n\in K\}\).  From \(k\in K\), take one ordinary
jump.  There are finitely many channel successors, and from each successor
the expected hitting time of \(K\) is finite by \eqref{eq:episodes-to-K}.
Thus \(\mathbb E_kT_K^+<\infty\) for every \(k\in K\).

Record successive visits to \(K\).  The resulting trace chain is finite and
irreducible.  Indeed, any augmented path between two points of \(K\), with
its intermediate \(K\)-visits retained, gives a positive-probability trace
path.  Fix \(k_*\in K\).  From every trace state choose a finite
positive-probability trace path to \(k_*\).  Finiteness supplies common
constants \(m<\infty\) and \(\varepsilon>0\) such that, from every trace
state, \(k_*\) is hit within the next \(m\) trace transitions with probability
at least \(\varepsilon\).  Starting from \(k_*\), first take one trace
transition and then apply this estimate in successive blocks of \(m\)
transitions.  The number of blocks before the next visit to \(k_*\) is
geometrically dominated, so the positive trace return time has finite mean.

Finally, set
\[
  B=\max_{k\in K}\mathbb E_kT_K^+<\infty.
\]
If \(M\) is the positive trace return count and \(L_j\) is the number of
original jumps in its \(j\)-th trace excursion, let \(\mathcal F_j\) be the
sigma-field at the start of that excursion.  Continue the trace chain after
its return to \(k_*\), so these objects are defined for all \(j\ge0\).  Then
\(\{j<M\}\in\mathcal F_j\), and the strong Markov property gives
\(\mathbb E[L_j\mid\mathcal F_j]\le B\).  Tonelli's theorem and the tail-sum
formula therefore yield
\[
\begin{aligned}
  \mathbb E\!\left[\sum_{j=0}^{M-1}L_j\right]
  &=\sum_{j\ge0}\mathbb E\!\left[
      \mathbf1_{\{j<M\}}L_j\right]\\
  &\le B\sum_{j\ge0}\mathbb P(j<M)
   =B\,\mathbb E M<\infty.
\end{aligned}
\]

\noindent
Hence the augmented embedded chain has finite expected positive return time
to \(k_*\).
\end{proof}

Projection away from the target mark shows that the population embedded chain
has finite expected positive return to the population coordinate \(x_*\) of
\(k_*\): its first positive population return occurs no later than the marked
return to \(k_*\).
It remains to construct physical time and rule out explosion.

\begin{proposition}[Embedded-chain-to-CTMC conversion]\label{prop:ctmc}
Let a countable irreducible embedded chain be positive recurrent.  Suppose a
pure-jump construction assigns a finite rate \(\Lambda(x)>0\) at every state
and that \(\inf_x\Lambda(x)>0\).  Then the resulting minimal CTMC is
nonexplosive and positive recurrent.
\end{proposition}

\begin{proof}
Let \(N\) be a finite-mean positive return count to a recurrent embedded state
\(x_*\), and let \(H_j\) be the holding times.  Conditional on the embedded
path, \(H_j\) is exponential with mean \(1/\Lambda(X_j)\).  Tonelli's theorem
and conditional expectation give
\[
 \mathbb E\sum_{j=0}^{N-1}H_j
 =\mathbb E\sum_{j\ge0}\mathbf1_{\{j<N\}}
   \frac1{\Lambda(X_j)}
 \le\frac{\mathbb EN}{\inf_x\Lambda(x)}<\infty.
\]
Thus the expected physical return time is finite, provided the construction
is nonexplosive.

Finite mean return makes \(x_*\) recurrent for the embedded chain, so it is
visited infinitely often almost surely.  On each visit, the following
holding time has the same exponential law of finite rate \(\Lambda(x_*)\).
The strong Markov construction makes these recurrent holding times
independent copies.  Their infinite sum diverges almost surely.  Since total
physical time dominates this subseries, jump times cannot accumulate at a
finite time.  The minimal CTMC is nonexplosive and has finite mean positive
return time to \(x_*\).
\end{proof}

Apply Proposition~\ref{prop:ctmc} directly to the population embedded chain on
\(\Gamma\).  Let \(\kappa_{\min}=\min_r\kappa_r>0\), with the minimum taken
over the finite reduced set of genuine channels.  Every \(x\in\Gamma\) enables
at least one such channel \(r\); otherwise \(x\) would be absorbing and the
infinite irreducible class \(\Gamma\) would be a singleton.  Since a positive
falling factorial is an integer,
\begin{equation}\label{eq:rate-lower}
  \Lambda(x)\ge\kappa_r(x)_{y_r}\ge\kappa_{\min},
  \qquad x\in\Gamma.
\end{equation}
Propositions~\ref{prop:trace} and \ref{prop:ctmc} therefore prove
Theorem~\ref{thm:main}.  For context, nonexplosion was previously obtained for
complex-balanced stochastic systems through a stationary-integrability
argument \citep{AndersonCappellettiKoyamaKurtz2018}.  For the present subclass,
the return-cycle construction recovers nonexplosion directly, consistently
with the more general second-order endotactic nonexplosion theorem of
\citet{Xu2026}.

The same return cycle yields Corollary~\ref{cor:stationary}.  Explicitly, for
the positive return time \(T_{x_*}^+\) defined in Section~\ref{sec:model},
\begin{equation}\label{eq:occupation-formula}
 \pi_\Gamma(y)=
 \frac{\displaystyle
   \mathbb E_{x_*}\!\left[\int_0^{T_{x_*}^+}
      \mathbf1_{\{X(t)=y\}}\,dt\right]}
 {\mathbb E_{x_*}T_{x_*}^+},
 \qquad y\in\Gamma,
\end{equation}
is the normalized mean occupation measure of one regenerative cycle.  The
standard regenerative occupation theorem makes this measure stationary
\citep{Asmussen2003}; its denominator normalizes it to one, and irreducibility
gives uniqueness.

\section{Scope and further questions}\label{sec:scope}

The proof is class-wise.  Coordinate faces, siphons, parity restrictions, and
other lattice constraints require no separate interior argument: every
constructed target-following path consists of actual enabled channels and
therefore remains in the fixed closed class.  Species absent from all
complexes are constant on each class and appear automatically in the
invariant branch of Lemma~\ref{lem:top}.

Three stability statements should be distinguished.  Theorem~\ref{thm:main}
proves positive recurrence on every communicating class that is already
closed.  It does not prove the additional Anderson--Kim stability property of
finite expected entrance, from an arbitrary initial state, into the union of
closed irreducible components.  The larger questions of multiple linkage
classes, molecularity above two, and the full weakly reversible stochastic
positive-recurrence conjecture also remain open.

The one-linkage hypothesis is used at one decisive point: for every carried
target \(t\) and selected terminal complex \(c\), a directed path
\eqref{eq:path} must exist.  With several linkage classes, a carried target
need not reach a terminal source belonging to another class, so the finite
path library does not directly glue.  That extension remains open here.

Bimolecularity is used in the exhaustive case split of
Lemma~\ref{lem:top}: a top complex either contains two particles supported on
divergent coordinates or exactly one such particle, leaving at most one
bounded companion species.  Complexes of molecularity three or higher admit
additional top configurations not covered by the argument.  No claim is made
for them.

The potentially reusable part of the proof is the combination of marked
reaction-channel augmentation, subtraction of the carried target, the exact
residual-factorial increment, finite target-following propagation, and
normalized-log top-complex compactification.  The multiple-linkage obstruction
is precise: a target in one linkage class need not have a directed path to a
useful terminal complex in another linkage class.  Overcoming that obstruction
would require an additional mechanism and is not a routine extension of the
present path argument.

The proof establishes positive recurrence and the existence of a stationary
probability distribution.  It does not provide a product form, explicit tail
bounds, exponential ergodicity, or quantitative mixing rates.  The finite set
\(K\) is obtained by analytic compactness, not by enumeration, and the proof
gives no practically useful universal estimate of its location or diameter.
Its scale can depend arbitrarily strongly on positive rate ratios, as the
exact example following Proposition~\ref{prop:K} illustrates.  Effective
rate-dependent recurrence, tail, and mixing estimates remain open.

\subsection*{Relevance to stochastic systems biology}

Chemical-master-equation models represent molecule-count noise in
low-copy-number biochemical networks.  When a model satisfies weak
reversibility, one linkage class, and bimolecularity, the theorem supplies a
parameter-independent structural conclusion: for every choice of positive
rate constants, each closed communicating class has one stationary
probability law.  Within such a class, the long-run molecule-count
distribution, state frequencies, and stationary averages of bounded
observables are therefore defined, and the process cannot escape permanently
to infinity.  Means, correlations, and other moments of unbounded count
observables are available only when they are integrable under that law.

The hypotheses are structural restrictions, not claims about all biological
networks.  The conclusion provides neither an explicit stationary formula nor
a mixing-time estimate, and it does not say that individual sample paths stay
inside a bounded region.  Large transient excursions remain possible, with no
quantitative bound supplied here.  The result is thus a qualitative stability
criterion for those intracellular, signalling, enzymatic, gene-regulatory, or
synthetic biochemical models whose effective reaction graphs fall within the
stated class; no empirical or organism-specific validation is asserted.

\section*{Declarations}

\paragraph{Author information.}
Alec Kriebel, Independent researcher.\\
Correspondence: \mbox{\href{mailto:me@aleckriebel.com}{me@aleckriebel.com}}.\\
ORCID: \href{https://orcid.org/0009-0001-9320-500X}{0009-0001-9320-500X}.

\paragraph{Code availability.}
The self-contained verification package and supporting materials are
available in the public project repository at
\url{https://github.com/AlecKriebel/Math/tree/bimolecular-positive-recurrence-v1.0/bimolecular_positive_recurrence_publication_v1}.
The submitted archive is also supplied with the manuscript so that its exact
Git snapshot and checksums can be retained with the submission record.  No
computation is used in the universal proof.

\paragraph{Preprint status.}
This is an unrefereed preprint prepared for expert audit and possible journal
or preprint-server submission.

\paragraph{Funding.}
This research received no specific grant from any funding agency in the public,
commercial, or not-for-profit sectors.

\paragraph{Competing interests.}
The author declares no competing interests.

\subsection*{Declaration of generative-AI use}
Generative-AI systems were used substantively for mathematical exploration,
counterexample attempts, proof criticism, verification code, literature
support, drafting, and editorial revision.  OpenAI ChatGPT, GPT-5.6 Pro,
accessed through the ChatGPT web interface on 5--6 August 2026, was used across
those activities.  Anthropic Claude, Opus 5 with Max effort, accessed through
claude.ai on 6 August 2026, was used for adversarial proof and manuscript
review; Claude Code with Opus 5, accessed through its command-line application
that day, was used for related publication formatting.  OpenAI Codex desktop,
using a GPT-5-family coding/research model whose exact backend identifier was
not exposed, was used on 8--9 August 2026 for publication-readiness revision,
audit, code work, and validation.  Rejected AI-assisted approaches are not
part of the final proof; the author directed the research, determined the
released claims, and assumes responsibility for the manuscript and artifacts.
No independent expert human validation is claimed, and no AI system is listed
as an author; a full tool-by-tool statement accompanies the paper.

\appendix
\section{Technical details for the trace chain and physical time}\label{app:trace}

This appendix records several interfaces that can otherwise be obscured by
standard Markov-chain shorthand.

First, the augmented state records the actual marked channel.  If two
channels have the same population displacement but different sources or
targets, their post-jump populations can agree while their carried targets do
not.  They therefore remain separate.  Only channels with identical source
and target may be aggregated.

Second, \(K\) in \eqref{eq:K} is nonempty, not merely finite.  The global
minimum argument in Proposition~\ref{prop:K} is pathwise: every possible
episode endpoint has potential at least the starting minimum.  No optional
stopping or integrability assertion is needed for this step because each
episode has at most \(|\mathcal C|\) jumps.

Third, the trace-chain proof distinguishes hitting from return.  The first
ordinary jump from \(k\in K\) forces a positive return time.  The finite
expected time back to \(K\) follows by conditioning on finitely many
successors.  The subsequent geometric-block argument is applied to the
finite trace chain, and only afterward are the random trace excursions
expanded into original embedded jumps.

Fourth, the embedded chain is constructed first.  Its recurrent population
state \(x_*\) is visited infinitely often.  The holding times attached to
those visits form
an infinite sequence of independent exponentials with one fixed finite rate,
whose sum diverges almost surely.  This directly prevents finite-time
accumulation.

Finally, singleton absorbing classes are handled before augmentation.  In an
infinite irreducible class every population state has an enabled genuine
channel; otherwise it would be absorbing and its class would be a singleton.

\section{Supplementary computational checks}\label{app:verification}

The accompanying dependency-free Python package performs exact-arithmetic
checks of the falling-factorial identity and entropy rewrite, verifies the
scalar-envelope branch conditions, and tests the top-complex classification
on a deterministic finite atlas with separately validated witnesses.  It also
contains fixed-seed falsification tests and exact finite calibrations of the
rate-degeneration example and stationary return-cycle formula.  No finite
atlas or random test proves the universal theorem, and simulations do not
prove recurrence.  In particular, the finiteness of \(K\) follows only from
the analytic compactness proof: the package neither enumerates \(K\) nor
certifies a practical upper bound on its size, diameter, or location.

\bibliographystyle{plainnat}
\bibliography{references}
